\section{Relative homotopy-coherent primitives on a coefficient tail}
\label{sec:136-derived-tail-primitives}

\paragraph{Archive and scope.}
This is a positive counterpart to ARCH-136-NATIVE-TAIL-NOSPLIT.
The relevant hypotheses are the two uniform admissibility conditions
of cote 136, physical pp.~117 and 119, footers 116 and 118,
scans 117 and 119; the common-source problem is on physical p.~121,
footer 120, scan 121. We prove a linear derived comparison and a
relative coherent Moore-pair lift. We do not infer that the resulting
nonflat parameter presentation gives an admissible strict PD extension.

Work over \(S=\Gamma_{\mathbb Z}(T)\). For \(M\geq1\), let
\(I_M=\bigoplus_{w\geq M}\mathbb ZT^{[w]}\).
Choose a weight-graded free \(S\)-resolution
\[
 \cdots\longrightarrow P_2\xrightarrow{\partial_2}P_1
       \xrightarrow{\partial_1}P_0
       \xrightarrow{\epsilon}I_M\longrightarrow0
                                                        \tag{136.DT1}
\]
in which every homogeneous free generator of every \(P_i\)
has weight at least \(M\). Such a resolution exists: choose
a free cover on homogeneous generators, then repeat on its
kernel. Every successive kernel is still concentrated in weights
at least \(M\), because the coefficient algebra has no negative
weights. No finite generation is required.
The complex \(P_\bullet\) is homological; its differential
lowers the index.

For a weight offset \(v\), write \(P_i[v]\) for the same
module with all weights increased by \(v\). If \(C\) is an
admissible theory, define a nonnegative cochain complex of
simplicial abelian groups by
\[
 \mathcal H_{P,v}(C)^n
       =\bigoplus_{i=0}^{n}
          \operatorname{Hom}^{0}_{S,\mathrm{gr}}
                  (P_i[v],C^{*,n-i}).                  \tag{136.DT2}
\]
Each Hom is a product of form rows of \(C\), one for each
homogeneous free generator. The direct sum over \(i\) is finite.
For a total degree \(n\) cochain \(f=(f_i)\), the differential is
\[
 (\delta f)_i=d_Cf_i-(-1)^n f_{i-1}\partial_i,
                 \qquad f_{-1}=0.                     \tag{136.DT3}
\]
It squares to zero by \(d_C^2=0\), \(\partial^2=0\), and
cancellation of the two mixed terms.
Let
\[
 K_{P,v}=\operatorname{Hom}^{0}_{S,\mathrm{gr}}
                         (P_\bullet[v],S[0])           \tag{136.DT4}
\]
with the same total convention. It is a nonnegative coefficient
cochain complex, usually not concentrated in degree zero.

\paragraph{ARCH-136-DERIVED-TAIL-COMPARISON.}
Fix \(d\geq0\). There is a number \(W_C(d)\), depending only
on the vertical and horizontal admissibility bounds of \(C\)
through form degree \(d+1\), such that if
\[
                         M+v\geq W_C(d),               \tag{136.DT5}
\]
then, for every simplicial pair \(Y\subseteq X\), the canonical
derived-section comparison induces natural isomorphisms
\[
 H^j\operatorname{Tot}\operatorname{Hom}^{0}_{S,\mathrm{gr}}
                 (P_\bullet[v],C(X,Y))
       \ \cong\ H^j(X,Y;K_{P,v})\qquad(0\leq j\leq d).
                                                        \tag{136.DT6}
\]
The right-hand side is ordinary relative cohomology with the
coefficient complex \(K_{P,v}\), namely the cohomology of
\(\operatorname{Tot}\operatorname{Hom}_{\mathbb Z}
(N(\mathbb Z[X]/\mathbb Z[Y]),K_{P,v})\).
In particular the left side is invariant under homotopy equivalences
of pairs in this range. The same bound works for arbitrary ranks
and lengths of the resolution (136.DT1).

\begin{proof}
Take \(W_C(d)\) to dominate all the stated vertical and horizontal
bounds through form degree \(d+1\), and to be positive.
Every component of (136.DT2) relevant to total degrees at most
\(d+1\) is a product of rows \(C^{w,q}\) with
\(w\geq M+v\) and \(q\leq d+1\).
Products of abelian groups are exact, and normalization of a
simplicial module commutes with products. Thus each of these
total form rows is horizontally acyclic.

There is a natural coefficient map
\[
             \underline{K_{P,v}}\longrightarrow
                              \mathcal H_{P,v}(C),      \tag{136.DT7}
\]
where the underline means constant in the original simplicial
direction. At every simplex it is a cohomology isomorphism
through degree \(d+1\). To see this, filter (136.DT2) by the
resolution index. In each fixed column, taking vertical cohomology
commutes with the product expressing Hom from the free module
\(P_i[v]\). The admissibility condition gives just the specified
coefficient copy in form zero and zero in positive forms through
the required range. The first-quadrant spectral sequence, whose
diagonals in total degrees at most \(d+1\) are finite, therefore
identifies this range with (136.DT4).

Here is the derived-section argument with its exact projectivity
requirement. Put \(L_{X,Y}=\mathbb Z[X]/\mathbb Z[Y]\).
Its normalization is a nonnegative homological complex of free
abelian groups, with the relative nondegenerate simplices as basis.
Use the projective resolution in the category of simplicial
abelian groups constructed from
\[
 Q(A)=S_0(A_0)\oplus\bigoplus_{r\geq1}D_r(A_r)
                       \longrightarrow A
                                                        \tag{136.DT8}
\]
on normalized complexes. Its kernel is the shifted tail
\(A'_r=A_{r+1}\), with differential \(-d_A\).
ARCH-136-TWISTED-COEFFICIENT-COMPARISON proved directly that
\(\operatorname{Ext}^{s}_{s\mathrm{Ab}}(L_{X,Y},B)=0\)
for \(s>0\) whenever \(B\) is a horizontally acyclic row:
a map from the shifted tail extends over (136.DT8) by recursively
lifting cycles, using freeness of each \(A_r\).
It did not require projectivity of the row \(B\).

Apply this resolution to (136.DT7). The row acyclicity proved
above compares ordinary sections with derived sections through
degree \(d\). Exactness of Hom from each projective term and
the simplexwise cohomology comparison in (136.DT7) compare those
derived sections with the coefficient complex through degree
\(d\). All the resulting double and triple complexes have
nonnegative cochain indices and finite diagonals in each fixed
total degree; the possibly infinite products are inside individual
terms and are exact. There is no exchange of scalar localization
with a section limit.

For a constant coefficient module the resolution (136.DT8)
computes the ordinary cochain complex on \(N(L_{X,Y})\).
Applying this termwise to the nonnegative complex \(K_{P,v}\)
gives precisely the right side of (136.DT6), again with finite
diagonals. Products commute with all section limits, so the ordinary
sections on the left are exactly the displayed total Hom into
\(C(X,Y)\). This proves (136.DT6), naturally in pairs and in
maps of coefficient resolutions.

A homotopy equivalence of simplicial pairs induces a chain-homotopy
equivalence of their relative normalized free chains. Precomposition
therefore induces a cochain-homotopy equivalence with coefficients
in \(K_{P,v}\). This proves the asserted homotopy invariance.
\end{proof}

\paragraph{ARCH-136-RELATIVE-COHERENT-MOORE.}
Let \(q\geq1\) be odd, let \(p\geq2\) be an integer, and put
\[
                         D=p(q+1),\qquad v=pm.         \tag{136.DT9}
\]
Suppose the input row \(C_*^{m,q}\) is horizontally acyclic,
including in degree zero. Take any simplicial pair \(Y\subseteq X\)
and any \(a\in C^{m,q}(X,Y)\). Write
\[
 b=da,\qquad z=\gamma_p(b),\qquad
 e=a\gamma_{p-1}(b),\qquad dz=0,\quad de=pz.          \tag{136.DT10}
\]
Choose \(M\) with \(M+v\geq W_C(D)\), and choose (136.DT1).
There are homogeneous \(S\)-linear maps
\[
 \begin{aligned}
 U_i &: P_i[v]\longrightarrow C^{*,D-1-i}(X,Y)
                              &&(0\leq i\leq D-1),\\
 V_i &: P_i[v]\longrightarrow C^{*,D-2-i}(X,Y)
                              &&(0\leq i\leq D-2),
 \end{aligned}                                         \tag{136.DT11}
\]
such that, writing \(f_z(u)=\epsilon(u)z\) and
\(f_e(u)=\epsilon(u)e\) on \(P_0[v]\), and zero in other
resolution indices,
\[
                         \delta U=f_z,\qquad
                         \delta V=pU-f_e.              \tag{136.DT12}
\]
This is a full homotopy-coherent coefficient-tail Moore pair.
All components vanish on \(Y\).

In particular the last component factors through the finite
resolution truncation
\[
 \overline P_{D-1}
        =\operatorname{coker}(P_D\longrightarrow P_{D-1}),
 \qquad U_{D-1}:\overline P_{D-1}[v]
                          \longrightarrow C^{*,0}(X,Y).
                                                        \tag{136.DT13}
\]
The last module is not asserted projective or flat.
The finite exact resolution ending with \(\overline P_{D-1}\)
and the maps (136.DT11) are actual algebraic data, not an infinite
formal sum. For finitely many targets, one may use the same \(M\),
the same resolution, and the same finite parameter presentation.
If \(Y\) contains every vertex, all their degree-zero parameter
values are zero at vertices.

\begin{proof}
First extend \(a\) and zero to an interpolation
\[
  \widetilde a\in C^{m,q}(X\times\Delta[1],Y\times\Delta[1])
                                                        \tag{136.DT14}
\]
which is \(a\) at one endpoint and zero at the other.
Prescribe these values on
\((X\times\partial\Delta[1])\cup(Y\times\Delta[1])\).
They agree on the intersections because \(a|_Y=0\).
The normalized relative chains of the remaining inclusion are
degreewise free. Extend the associated map into the normalized
input row one relative generator at a time; its already assigned
boundary is a cycle, and acyclicity of that row supplies a lift.
This is the same row-extension argument as
ARCH-136-GUARDED-PAIR, now with the additional relative subcomplex.
It requires no chosen linear contraction.

Let \(\mathcal A=\operatorname{Tot}
\operatorname{Hom}^{0}_{S,\mathrm{gr}}
(P_\bullet[v],C(X,Y))\), and for any cochain complex define
\[
 P_p(\mathcal A)^r=\mathcal A^r\oplus\mathcal A^{r+1},
 \qquad d(x,y)=(\delta x-py,-\delta y).
                                                        \tag{136.DT15}
\]
It has the natural exact sequence
\[
 0\longrightarrow H^r(\mathcal A)/pH^r(\mathcal A)
 \longrightarrow H^r(P_p(\mathcal A))
 \longrightarrow H^{r+1}(\mathcal A)[p]
 \longrightarrow0.                                    \tag{136.DT16}
\]
No flatness is needed for this mapping-cone sequence.
Equations (136.DT10) and \(\epsilon\partial_1=0\) make
\((f_e,f_z)\) a closed element of (136.DT15) in degree
\(D-1\).

Apply the same construction to \(\widetilde a\). It gives a
closed relative cone element over
\((X\times\Delta[1],Y\times\Delta[1])\), with endpoints
\((f_e,f_z)\) and zero. The projection of this pair to \((X,Y)\)
is a homotopy equivalence. By ARCH-136-DERIVED-TAIL-COMPARISON
it induces cohomology isomorphisms on \(\mathcal A\) through
degree \(D\). Equation (136.DT16) then gives an isomorphism
on cone cohomology in degree \(D-1\). Both endpoint maps are
left inverses of this isomorphism, so they are equal there.
Consequently \((f_e,f_z)\) has zero cone cohomology class.

Choose a cone primitive \((r,t)\), of degree \(D-2\), with
\[
             (f_e,f_z)=(\delta r-pt,-\delta t).
\]
Put \(U=-t\), \(V=-r\). These have total degrees \(D-1\)
and \(D-2\), respectively, and satisfy (136.DT12).
Because the original form grading is nonnegative, their only
possible components are exactly (136.DT11). Every component
lies in the relative complex and therefore vanishes on \(Y\).

For clarity, \(D\) is even. Formula (136.DT3) gives the full
component relations, with omitted indices interpreted as zero:
\[
 \begin{split}
 d_C U_i+U_{i-1}\partial_i&=(f_z)_i,\\
 d_C V_i-V_{i-1}\partial_i&=pU_i-(f_e)_i.
 \end{split}                                           \tag{136.DT17}
\]
At index \(D\) the first relation is
\(U_{D-1}\partial_D=0\), proving (136.DT13).
At index \(D-1\) the second relation is
\[
                -V_{D-2}\partial_{D-1}=pU_{D-1}.
                                                        \tag{136.DT18}
\]
These include the bottom relation of the coherent Moore pair;
it is not silently discarded when the resolution is truncated.
Exactness of (136.DT1) makes the finite sequence
\[
 0\longrightarrow\overline P_{D-1}\longrightarrow P_{D-2}
       \longrightarrow\cdots\longrightarrow P_0
       \longrightarrow I_M\longrightarrow0
                                                        \tag{136.DT19}
\]
exact. This is the stated bounded, possibly nonflat, parameter
resolution.

For a finite list of target theories, take the maximum of their
cutoffs in (136.DT5), then use one resolution of that common tail.
Each interpolation and each cone primitive may be chosen separately;
their maps nevertheless satisfy precisely the same linear parameter
relations (136.DT17)--(136.DT18). When \(Y\) contains the vertices,
all their bottom values are the same coefficient value, namely zero.
This proves every assertion.
\end{proof}

\paragraph{Exact consequence and remaining gate.}
The direct coefficient-linear primitive prohibited by
ARCH-136-NATIVE-TAIL-NOSPLIT is replaced by a full coherent lift
on a projective resolution, terminating in an actual nonflat
syzygy module in form zero. Relative interpolation makes its
vertex values simultaneously zero in every target. Thus neither
the absence of a direct tail splitting nor the infinite ordinary
projective dimension of a tail by itself prevents these finite
coherent native lifts.

This statement still reaches form degree zero for every \(p\).
It provides no lower form cutoff tending to infinity with \(p\).
The last syzygy is generally nonflat, and free PD attachment of
this parameter presentation may change both vertical cohomology
and horizontal row bounds. Nor is every new cycle after such an
attachment a universal derivative cycle of the form (136.DT10).
Strict PD attachment, preservation of the two admissibility
conditions, and iteration beyond these native operations remain
separate, unresolved construction requirements.
